\section{Method: $p$-fair BPE}
\label{sec:method}

Like classical BPE, $p$-fair BPE builds the vocabulary one merge at a time and chooses each
merge greedily: we derive the merge rule implied by Eq.~\eqref{eq:obj}, show that it contains
classical and Parity-aware BPE as limits, and give an exact and fast training algorithm.

\subsection{The greedy merge rule}
\label{sec:rule}

Merging a pair $(a,b)$ replaces each of its occurrences by one symbol, so the token count
of a text falls by the number of occurrences of the pair. The pair counts $c_\ell$ come from
the training text, whereas the costs $t_\ell$ are measured on the development set. To
connect the two, we assume that a pair occurs in the development text of language $\ell$ at
the same rate as in its training text. Let $N_\ell$ denote the size of $\mathcal{D}_\ell$
expressed in development sentences, that is, its length divided by the average length of a
development sentence in language $\ell$. Under this assumption, merging $(a,b)$ lowers
$T_\ell$ by $S\,c_\ell(a,b)/N_\ell$ and therefore lowers $t_\ell$ by $c_\ell(a,b)/N_\ell$.

\begin{proposition}
\label{prop:rule}
Let $\pi_\ell = N_\ell/\sum_k N_k$. To first order, the merge that decreases $F_p$ the most
is
\begin{equation}
  (a,b)^\star \;=\; \argmax_{(a,b)}\; s_p(a,b), \qquad
  s_p(a,b) = \sum_{\ell\in\mathcal{L}} t_\ell^{\,p-1}\, c_\ell(a,b).
  \label{eq:rule}
\end{equation}
\end{proposition}
The proof, a first-order expansion of $F_p$, is in Appendix~\ref{app:proofs}.


Thus, the rule is a pair count in which each language is weighted by its current cost raised
to the power $p-1$. The weights change after every merge; we update every $t_\ell$ exactly
by applying the merge to the development set, so only the per-merge choice relies on the
first-order approximation.

With Hindi at 35.0 tokens per sentence and English at 29.0, as under classical BPE in our
experiments, a Hindi pair counts $(35.0/29.0)^{p-1}$ times as much as an English pair of the
same frequency: 1.2 times at $p=2$ and 3.7 times at $p=8$.

\subsection{Classical and Parity-aware BPE as limits}
\label{sec:limits}

The rule in Eq.~\eqref{eq:rule} has two limits. At $p=1$ every weight equals one,
$s_1(a,b)=\sum_\ell c_\ell(a,b)$ is the count of the pair in the pooled training text, and
the rule is exactly classical BPE. As $p\to\infty$, dividing $s_p$ by
$(\max_\ell t_\ell)^{p-1}$ does not change the $\argmax$; the weight of the worst-off
language is then one and every other weight tends to zero, so the rule selects the most
frequent pair of the worst-off language. This is the base rule of Parity-aware BPE
\citep{foroutan2026} with cost measured per parallel sentence, and corresponds to a
subgradient step on $\max_\ell t_\ell$. For $1<p<\infty$ every language keeps a positive weight: unlike the parity rule (one language
per merge) and the Worst-$N$ rule \citep{kanjirangat2026} (the $N$ worst-off weighted
equally), $p$-fair BPE weights all languages continuously by their cost, and unlike the
hybrid schedule \citep{foroutan2026} it uses one parameter that is constant throughout
training.

\subsection{Exact training by lazy evaluation}
\label{sec:lazy}

Efficient implementations of classical BPE keep the pair counts in a priority queue and
update only the pairs that a merge touches. In Eq.~\eqref{eq:rule} the weights change after
every merge, so every score changes, and rescoring all $|\mathcal{P}|$ pairs at each of the
$\sim$31{,}000 steps costs $O(|\mathcal{P}|\,L)$ per step. The following property avoids this.

\begin{proposition}
\label{prop:lazy}
For $p\ge1$, the score $s_p(a,b)$ of a pair that already exists never increases during
training.
\end{proposition}
The proof is in Appendix~\ref{app:proofs}.


Consequently, a stored score is always an upper bound on the current score. We keep the
pairs in a max-heap keyed by their stored scores, pop the top entry and recompute its score
with the current weights: if it is unchanged, the entry is the exact maximiser of
Eq.~\eqref{eq:rule}, since every other stored score bounds its current score from above;
otherwise we push it back with the new score and repeat (Algorithm~\ref{alg:pfair} in
Appendix~\ref{app:impl}). This is the lazy greedy technique of \citet{minoux1978}, and it
returns the same merge as full rescoring at every step. For $p=\infty$ we keep one heap per
language keyed by $c_\ell$ and pop from the heap of the worst-off language.
